\section{Related Work}

\subsection{Persistent-memory poisoning as a threat class}

OWASP's Agentic Security Initiative names cross-session memory
poisoning as a distinct threat category (Memory Poisoning, threat ID
T1 in its Agentic AI Threats and Mitigations guide) \cite{owasp-asi},
separate from single-turn prompt injection precisely because the
compromised state persists and can influence multiple future
interactions. A cluster of recent work---MemLineage, MemAudit,
MemSecBench, MPBench, AgentPoison, and MINJA---has established this as
an active subfield within months. We position our contribution against
all six.

\subsection{Preventive lineage enforcement: MemLineage}

MemLineage (arXiv 2605.14421) \cite{memlineage} builds a
lineage-guided enforcement system: a cryptographic chain-of-custody
layer plus an LLM-mediated derivation-lineage DAG (module M4). M4
records an edge $(p \to c)$ whenever parent memory $p$ was present in
the LLM's context when child memory $c$ was written---execution/derivation
provenance, not lexical or cryptographic matching, and the paper states
this is specifically designed to survive paraphrasing and
summarization. Three attribution mechanisms weight parent$\to$child
edges (Coarse: uniform weight 1.0, maximum recall/minimum precision;
LmSelfEval: LLM-judged semantic influence; AttnAttr: white-box
attention attribution), and the guarantee that an untrusted label
reaches a downstream chain tip is \emph{conditional} on every edge on
the path exceeding an attribution threshold $\tau$---MemLineage's own
$\tau \times K$ sweep and $r_K$ metric (strong-edge recall at chain
length $K$) study exactly this degradation. Their evaluation is a
deterministic scripted-judge harness plus live runs on two models,
evaluated primarily over \textbf{linear} derivation chains, answering:
does a critical ancestry path remain detectable, so that a
sensitive-action gate can catch it, as attribution threshold and depth
vary?

Our work sits downstream of a different trigger condition. MemLineage
asks whether an untrusted label survives \emph{forward} to gate a
\emph{future} action, before harm occurs. We ask: given a memory that
is \emph{already known} to be compromised, what is the actual shape of
its downstream footprint in a \textbf{branching} graph (multiple
derived memories per retrieval event, not one), and what does
reconstructing and containing that footprint cost? Put plainly:
\textbf{prior lineage systems use provenance primarily as an
enforcement primitive; we treat provenance as an incident-response
primitive}, evaluated against a different object (reconstruction
accuracy and containment cost, not action-gating accuracy) and against
semantic downstream ground truth in branching derivation graphs rather
than linear chains --- branching structure alone is not the novelty
claim; the different trigger condition, ground truth, metrics, and
decision objective together are. MemLineage's Coarse
attribution mode acknowledges the precision/recall tradeoff we study
exists, but does not characterize it under branching write fan-out the
way we do (H1), and its evaluation does not construct or measure a
blast-radius precision/recall/inflation triple against ground-truth
semantic contamination the way our metrics do.

\subsection{Post-hoc attribution: MemAudit}

MemAudit (arXiv 2605.23723) \cite{memaudit} operates after a harmful
event has already been observed: given a defined harmful event $e =
(q^*, y^*, R^*)$, it combines a counterfactual memory-influence score
with a structural memory-consistency graph to rank suspicious memories
for targeted removal. Evaluated on MINJA QA and RAP/WebShop tasks
across three models, it reduces attack success rate to zero at
low-to-moderate contamination ratios, with a sharp operating-boundary
transition (QA: $\rho{=}0.20 \to$ 0\% ASR-after, $\rho{=}0.25 \to$
60--70\%; RAP: $\rho{=}0.15 \to$ 0\%, $\rho{=}0.24 \to$ 80--87\%) rather
than gradual degradation. Critically, MemAudit does \textbf{not}
maintain persistent derivation provenance connecting a known-compromised
source to downstream writes independent of replaying each harmful
event---its memory-consistency graph is a semantic \emph{consistency}
graph, not a derivation \emph{lineage} graph, and its own limitations
section states undetected failures cannot be repaired by the system.

We do not claim MemAudit cannot handle laundered content that is itself
later retrieved for a harmful event (its counterfactual replay and
anomaly detection may still catch such cases). The defensible
distinction is structural: MemAudit requires an observed harmful signal
to trigger investigation; we study reconstructing exposure from a
known-compromised \emph{source}, independent of whether any downstream
write has yet caused visible harm.

\subsection{Lifecycle benchmarking: MemSecBench}

MemSecBench (arXiv 2607.27080) \cite{memsecbench} is the closest match
in evaluation \emph{methodology}, though not in research question. It
traces a linked Write$\to$Execute$\to$Forget lifecycle anchored to a
case-specific \texttt{target\_memory\_manifest} describing malicious
semantics, evaluated across 310 linked test cases, 2 agent harnesses, 3
LLM backends, and 4 memory backends. It measures whether malicious
semantics persist (84.2\% configuration-macro-averaged persistence),
complete an end-to-end Write$\to$Execute chain (50.3\%), and can be
selectively repaired without damaging benign memory (56.1\% joint
success, vs.\ 86.3\% for target-removal alone)---the gap between those
last two figures is direct independent evidence that collateral damage
to benign memory, not detection, is repair's hard part, which motivates
our own \texttt{inflation\_ratio}/\texttt{containment\_cost} metrics.
MemSecBench does not construct a derivation-lineage graph, does not
track which \emph{other} memories become descendants of a compromised
one, and does not measure blast-radius precision/recall under branching
retrieval/write fan-out---it is semantic lifecycle tracking, not
structural provenance. We also adopt its case-construction precedent
(structured scenario authoring, human review gate, ground truth
anchored to a single semantic-target definition per case) for our own
scenario corpus.

\subsection{Attack taxonomy and detection-strength gradient: MPBench}

MPBench (arXiv 2606.04329v2 \cite{mpbench}---distinct from an unrelated
same-named paper, arXiv 2503.12505, which we do not cite) evaluates a
6-class attack taxonomy (Explicit Command Insertion, Conditional
Command Insertion, Salience-Driven Compaction Poisoning,
Policy-Conformant Fact Injection, False Precedent Insertion,
Skill-Procedure Insertion) across 3,240 adversarial and 2,997 benign
cases, two agents, seven domains. Our \texttt{poison\_form} axis is
\emph{not} this taxonomy---only \texttt{direct\_instruction} and
partially \texttt{embedded\_fact}/\texttt{authoritative\_framing}
overlap it; \texttt{multi\_hop\_setup} is our own construction.
MPBench's strong/weak-signal distinction (e.g.\ a prompt-injection
detector, PromptArmor, dropping from 84.44\% to 42.50\% detection
between strong- and weak-signal payloads) independently motivates our
H3 surface-vs-execution-provenance comparison. MPBench's
Skill-Procedure Insertion finding also shows downstream
dependency/amplification through repeated execution is already
discussed in prior work---we do not claim to be first to observe that
poison can compound downstream, only that no prior work constructs or
evaluates a structural derivation graph, reconstructs blast radius from
a known-compromised root, or measures over-tainting under branching
fan-out the way we do. MPBench itself recommends write-path provenance
tracking as a defense direction, so provenance tracking per se is not
our novelty claim.

\subsection{Attack-construction work: AgentPoison and MINJA}

AgentPoison (arXiv 2407.12784) \cite{agentpoison} and MINJA (arXiv
2503.03704) \cite{minja} are attack-construction papers on a different
object than ours. AgentPoison uses constrained-optimization trigger
generation to map triggered queries to a compact embedding region,
achieving ~82\% average poisoned-retrieval success (ASR-r) and ~63\%
end-to-end attack success (ASR-t) across three agent settings, but
assumes an attacker with white-box retriever access and the ability to
directly inject key-value pairs into the memory store---a stronger
capability assumption than ours. MINJA achieves ~98.2\% injection
success and ~76.8\% attack success by having the attacker behave as an
ordinary user (no direct memory access), via bridging steps and
progressive query shortening, but its main threat model assumes a
\emph{shared} memory bank across users, and its empirical results are
not validated for isolated per-user memory. Neither paper constructs or
evaluates a downstream derivation graph, blast-radius precision/recall,
or containment cost after compromise is known. We do borrow MINJA's
methodological observation that attack-instance content variation (they
use 9 independent victim-target pairs) matters more than run-level
stochasticity---independent support for our \texttt{scenario\_id}-based
sampling design (Section~\ref{sec:method}).

\subsection{Summary}

\begin{table}[h]
\centering
\caption{Positioning against prior memory-security work.}
\label{tab:related-work-summary}
\begin{tabular}{p{1.4cm}p{1.5cm}p{1.6cm}p{2.5cm}}
\toprule
 & Trigger & Graph shape & Question \\
\midrule
MemLineage & preventive (before action) & mostly linear chains & does ancestry survive to gate an action? \\
MemAudit & reactive (harm observed) & none (consistency graph) & which memory caused \emph{this} failure? \\
MemSecBench & lifecycle checkpoints & none (semantic manifest) & does malicious semantics persist/propagate/repair? \\
\textbf{This work} & \textbf{known compromise} & \textbf{branching derivation graph} & \textbf{what is exposed downstream, and what does containment cost?} \\
\bottomrule
\end{tabular}
\end{table}

We do not claim priority on persistent-memory poisoning as a threat, on
provenance tracking as a mitigation mechanism, or on the observation
that poison can compound downstream through repeated execution---all
three are established or independently suggested by the papers above.
The contribution is the specific research object: post-incident
blast-radius reconstruction in \textbf{branching} agent-memory graphs,
with metrics (blast-radius precision/recall, inflation ratio,
containment cost) that none of the six papers above measure.
